\chapter{Help a Lost Tourist in Your Town}

% ================= PERCAKAPAN 1-7 =================
\begin{convobox}{1}
    \noindent
    \jpkana{あ、}{A,} \jpkana{すみません。}{sumimasen.} 

    \vspace{1.5em}
    \noindent{\Large \color{articolor} \textbf{Arti:} Ah, permisi.}
\end{convobox}

\begin{convobox}{2}
    \noindent
    \jpword{いま}{今}{Ima}\jpkana{、}{,} 
    \jpkana{ちょっと}{chotto} 
    \jpkana{いいですか？}{ii desu ka?}

    \vspace{1.5em}
    \noindent{\Large \color{articolor} \textbf{Arti:} Apakah kamu ada waktu sebentar sekarang?}
\end{convobox}

\begin{convobox}{3}
    \noindent
    \jpword{びじゅつかん}{美術館}{Bijutsukan} \jpkana{に}{ni} 
    \jpword{い}{行}{i}\jpkana{きたいんですが、}{kitai n desu ga,}

    \vspace{1.5em}
    \noindent{\Large \color{articolor} \textbf{Arti:} Saya ingin pergi ke museum seni,}
\end{convobox}

\begin{convobox}{4}
    \noindent
    \jpkana{この}{Kono} 
    \jpword{ちか}{近}{chika}\jpkana{く}{ku} \jpkana{に}{ni} 
    \jpkana{ありますか？}{arimasu ka?}

    \vspace{1.5em}
    \noindent{\Large \color{articolor} \textbf{Arti:} Apakah ada di dekat sini?}
\end{convobox}

\begin{convobox}{5}
    \noindent
    \jpkana{はい。}{Hai.} \jpkana{ありますよ。}{Arimasu yo.}

    \vspace{1.5em}
    \noindent{\Large \color{articolor} \textbf{Arti:} Ya. Ada kok.}
\end{convobox}

\begin{convobox}{6}
    \noindent
    \jpword{ある}{歩}{Aru}\jpkana{いて}{ite} 
    \jpword{なんぷん}{何分}{nanpun} 
    \jpkana{ぐらいですか？}{gurai desu ka?}

    \vspace{1.5em}
    \noindent{\Large \color{articolor} \textbf{Arti:} Kira-kira berapa menit kalau jalan kaki?}
\end{convobox}

\begin{convobox}{7}
    \noindent
    \jpword{ある}{歩}{Aru}\jpkana{いて}{ite} 
    \jpword{にじゅっぷん}{20分}{ni-juppun} 
    \jpkana{ぐらいです。}{gurai desu.}

    \vspace{1.5em}
    \noindent{\Large \color{articolor} \textbf{Arti:} Kira-kira 20 menit jalan kaki.}
\end{convobox}

% % --- LESSON BUFFER 1 ---
\begin{lessonbox}{Catatan Belajar: Nanya Tempat Waktu}
    Wah, ada turis yang nyasar nih! Kalau kamu mau bertanya apakah ada suatu tempat di dekat sini, kamu bisa bilang: \textbf{"Kono chikaku ni arimasu ka?"} (Apakah ada di dekat sini?). \\
    
    Terus, kalau kamu mau tahu berapa lama perjalanannya kalau jalan kaki, gunakan \textbf{"Aruite nanpun gurai desu ka?"}. Yuk kita lihat apakah turisnya mau jalan kaki selama 20 menit!
\end{lessonbox}

% ================= PERCAKAPAN 8-16 =================
\begin{convobox}{8}
    \noindent
    \jpword{けっこう}{結構}{Kekkou} 
    \jpword{とお}{遠}{too}\jpkana{いですね。}{i desu ne.}

    \vspace{1.5em}
    \noindent{\Large \color{articolor} \textbf{Arti:} Lumayan jauh ya.}
\end{convobox}

\begin{convobox}{9}
    \noindent
    \jpkana{タクシー}{Takushii} \jpkana{は}{wa} 
    \jpword{たか}{高}{taka}\jpkana{いですか？}{i desu ka?}

    \vspace{1.5em}
    \noindent{\Large \color{articolor} \textbf{Arti:} Apakah taksi mahal?}
\end{convobox}

\begin{convobox}{10}
    \noindent
    \jpkana{タクシー}{Takushii} \jpkana{は}{wa} 
    \jpkana{ちょっと}{chotto} 
    \jpword{たか}{高}{taka}\jpkana{いと}{i to} 
    \jpword{おも}{思}{omo}\jpkana{います。}{imasu.}

    \vspace{1.5em}
    \noindent{\Large \color{articolor} \textbf{Arti:} Menurutku taksi sedikit mahal.}
\end{convobox}

\begin{convobox}{11}
    \noindent
    \jpword{でんしゃ}{電車}{Densha} \jpkana{とか}{toka} 
    \jpkana{バス}{basu} \jpkana{は}{wa} 
    \jpkana{ありますか？}{arimasu ka?}

    \vspace{1.5em}
    \noindent{\Large \color{articolor} \textbf{Arti:} Apakah ada kereta atau bus?}
\end{convobox}

\begin{convobox}{12}
    \noindent
    \jpkana{バス}{Basu} \jpkana{は}{wa} 
    \jpkana{ありませんが、}{arimasen ga,} 
    \jpword{でんしゃ}{電車}{densha} \jpkana{は}{wa} 
    \jpkana{ありますよ。}{arimasu yo.}

    \vspace{1.5em}
    \noindent{\Large \color{articolor} \textbf{Arti:} Bus tidak ada, tapi kereta ada lho.}
\end{convobox}

\begin{convobox}{13}
    \noindent
    \jpword{えき}{駅}{Eki} \jpkana{は}{wa} 
    \jpkana{どこ}{doko} \jpkana{に}{ni} 
    \jpkana{ありますか？}{arimasu ka?}

    \vspace{1.5em}
    \noindent{\Large \color{articolor} \textbf{Arti:} Stasiunnya ada di mana?}
\end{convobox}

\begin{convobox}{14}
    \noindent
    \jpkana{この}{Kono} 
    \jpword{みち}{道}{michi} \jpkana{を}{o} 
    \jpkana{まっすぐ}{massugu} 
    \jpword{い}{行}{i}\jpkana{って、}{tte,}

    \vspace{1.5em}
    \noindent{\Large \color{articolor} \textbf{Arti:} Ikuti jalan ini lurus terus,}
\end{convobox}

\begin{convobox}{15}
    \noindent
    \jpword{かど}{角}{Kado} \jpkana{を}{o} 
    \jpword{みぎ}{右}{migi} \jpkana{に}{ni} 
    \jpword{ま}{曲}{ma}\jpkana{がると、}{garu to,} 
    \jpword{えき}{駅}{eki} \jpkana{が}{ga} 
    \jpkana{あります。}{arimasu.}

    \vspace{1.5em}
    \noindent{\Large \color{articolor} \textbf{Arti:} Kalau belok kanan di tikungan, stasiunnya ada di sana.}
\end{convobox}

\begin{convobox}{16}
    \noindent
    \jpkana{ありがとうございます。}{Arigatou gozaimasu.}

    \vspace{1.5em}
    \noindent{\Large \color{articolor} \textbf{Arti:} Terima kasih banyak.}
\end{convobox}

% --- LESSON BUFFER 2 ---
\begin{lessonbox}{Catatan Belajar: Memberi Petunjuk Jalan!}
    Ternyata turisnya mau naik kereta! Saat kamu mau memberi tahu arah jalan ke seseorang, kamu bisa memakai kosakata ini:
    \begin{itemize}
        \item \textbf{Massugu} = Lurus terus
        \item \textbf{Migi} = Kanan (kalau Kiri = \textit{Hidari})
        \item \textbf{Magaru} = Belok
    \end{itemize}
    Jadi, \textit{"Massugu itte, migi ni magaru"} artinya "Lurus terus, lalu belok kanan". Keren banget kan bisa bantu turis nyasar!
\end{lessonbox}

% ================= PERCAKAPAN 17-22 =================
\begin{convobox}{17}
    \noindent
    \jpkana{この}{Kono} 
    \jpword{まち}{街}{machi} \jpkana{で}{de} 
    \jpword{ゆうめい}{有名}{yuumei}\jpkana{な}{na} 
    \jpword{かんこうち}{観光地}{kankouchi} \jpkana{や、}{ya,}

    \vspace{1.5em}
    \noindent{\Large \color{articolor} \textbf{Arti:} Tempat wisata yang terkenal di kota ini, atau}
\end{convobox}

\begin{convobox}{18}
    \noindent
    \jpkana{おすすめ}{Osusume} \jpkana{の}{no} 
    \jpword{ばしょ}{場所}{basho} \jpkana{は}{wa} 
    \jpkana{ありますか？}{arimasu ka?}

    \vspace{1.5em}
    \noindent{\Large \color{articolor} \textbf{Arti:} apakah ada tempat yang kamu rekomendasikan?}
\end{convobox}

\begin{convobox}{19}
    \noindent
    \jpword{おお}{大}{Oo}\jpkana{きな}{kina} 
    \jpword{ひろば}{広場}{hiroba} \jpkana{が}{ga} 
    \jpword{ゆうめい}{有名}{yuumei} \jpkana{です。}{desu.}

    \vspace{1.5em}
    \noindent{\Large \color{articolor} \textbf{Arti:} Alun-alun yang besar sangat terkenal lho.}
\end{convobox}

\begin{convobox}{20}
    \noindent
    \jpkana{いろいろな}{Iroiro na} 
    \jpword{みせ}{お店}{o-mise} \jpkana{が}{ga} \jpkana{あって}{atte} 
    \jpword{おも}{面}{omo}\jpword{しろ}{白}{shiro}\jpkana{いですよ。}{i desu yo.}

    \vspace{1.5em}
    \noindent{\Large \color{articolor} \textbf{Arti:} Ada bermacam-macam toko dan tempatnya sangat menarik.}
\end{convobox}

\begin{convobox}{21}
    \noindent
    \jpkana{ときどき}{Tokidoki} 
    \jpkana{イベント}{ibento} \jpkana{も}{mo} 
    \jpkana{あります。}{arimasu.}

    \vspace{1.5em}
    \noindent{\Large \color{articolor} \textbf{Arti:} Kadang-kadang juga ada acara (event).}
\end{convobox}

\begin{convobox}{22}
    \noindent
    \jpkana{わかりました。}{Wakarimashita.} 
    \jpkana{あとで}{Ato de} 
    \jpword{い}{行}{i}\jpkana{ってみます。}{tte mimasu.}

    \vspace{1.5em}
    \noindent{\Large \color{articolor} \textbf{Arti:} Saya mengerti. Nanti saya akan coba pergi ke sana.}
\end{convobox}

% --- LESSON BUFFER 3 ---
\begin{lessonbox}{Catatan Belajar: Rekomendasi Tempat!}
    Kalau kamu jalan-jalan dan bingung mau ke mana, kata ajaibnya adalah \textbf{おすすめ (Osusume)} yang berarti "Rekomendasi". \\
    
    Orang Jepang biasanya akan menjawab dengan antusias dan memberitahu tempat yang \textbf{有名 (Yuumei)} alias terkenal, atau yang \textbf{面白い (Omoshiroi)} alias menarik/seru!
\end{lessonbox}

% ================= PERCAKAPAN 23-34 =================
\begin{convobox}{23}
    \noindent
    \jpkana{そういえば、}{Sou ieba,} 
    \jpkana{のどが}{nodo ga} 
    \jpkana{かわいた}{kawaita} 
    \jpkana{ので}{node}

    \vspace{1.5em}
    \noindent{\Large \color{articolor} \textbf{Arti:} Ngomong-ngomong, karena saya merasa haus,}
\end{convobox}

\begin{convobox}{24}
    \noindent
    \jpword{みず}{水}{Mizu} \jpkana{を}{o} 
    \jpword{か}{買}{ka}\jpkana{いたいんですが、}{itai n desu ga,}

    \vspace{1.5em}
    \noindent{\Large \color{articolor} \textbf{Arti:} saya ingin membeli air,}
\end{convobox}

\begin{convobox}{25}
    \noindent
    \jpkana{この}{Kono} 
    \jpword{ちか}{近}{chika}\jpkana{く}{ku} \jpkana{に}{ni} 
    \jpkana{スーパー}{suupaa} \jpkana{とか}{toka} 
    \jpkana{コンビニ}{konbini} \jpkana{は}{wa} 
    \jpkana{ありますか？}{arimasu ka?}

    \vspace{1.5em}
    \noindent{\Large \color{articolor} \textbf{Arti:} apakah di dekat sini ada supermarket atau minimarket?}
\end{convobox}

\begin{convobox}{26}
    \noindent
    \jpkana{はい。}{Hai.} 
    \jpkana{あの}{ano} 
    \jpword{きいろ}{黄色}{kiiro}\jpkana{い}{i} 
    \jpword{かんばん}{看板}{kanban}\jpkana{、}{,} 
    \jpword{み}{見}{mi}\jpkana{えますか？}{emasu ka?}

    \vspace{1.5em}
    \noindent{\Large \color{articolor} \textbf{Arti:} Ya. Apakah kamu bisa melihat papan reklame kuning itu?}
\end{convobox}

\begin{convobox}{27}
    \noindent
    \jpkana{あれが}{Are ga} 
    \jpkana{コンビニ}{konbini} \jpkana{です。}{desu.}

    \vspace{1.5em}
    \noindent{\Large \color{articolor} \textbf{Arti:} Itu adalah minimarket.}
\end{convobox}

\begin{convobox}{28}
    \noindent
    \jpkana{わかりました。}{Wakarimashita.} 
    \jpkana{ありがとうございます。}{Arigatou gozaimasu.}

    \vspace{1.5em}
    \noindent{\Large \color{articolor} \textbf{Arti:} Saya mengerti. Terima kasih banyak.}
\end{convobox}

\begin{convobox}{29}
    \noindent
    \jpkana{あ、}{A,} 
    \jpkana{もうすぐ}{mou sugu} 
    \jpword{ひる}{お昼}{o-hiru} \jpkana{ですね。}{desu ne.}

    \vspace{1.5em}
    \noindent{\Large \color{articolor} \textbf{Arti:} Ah, sebentar lagi makan siang ya.}
\end{convobox}

\begin{convobox}{30}
    \noindent
    \jpkana{この}{Kono} 
    \jpword{ちか}{近}{chika}\jpkana{く}{ku} \jpkana{に}{ni} 
    \jpkana{おすすめの}{osusume no} 
    \jpkana{レストラン}{resutoran} \jpkana{や}{ya}

    \vspace{1.5em}
    \noindent{\Large \color{articolor} \textbf{Arti:} Di dekat sini, apakah ada restoran yang direkomendasikan, atau}
\end{convobox}

\begin{convobox}{31}
    \noindent
    \jpkana{カフェ}{kafe} \jpkana{は}{wa} 
    \jpkana{ありますか？}{arimasu ka?}

    \vspace{1.5em}
    \noindent{\Large \color{articolor} \textbf{Arti:} kafe?}
\end{convobox}

\begin{convobox}{32}
    \noindent
    \jpword{えき}{駅}{Eki} \jpkana{の}{no} 
    \jpword{ちか}{近}{chika}\jpkana{く}{ku} \jpkana{に}{ni} 
    \jpkana{『Pomodoro』}{Pomodoro} \jpkana{という}{to iu}

    \vspace{1.5em}
    \noindent{\Large \color{articolor} \textbf{Arti:} Di dekat stasiun ada tempat bernama 'Pomodoro', yaitu}
\end{convobox}

\begin{convobox}{33}
    \noindent
    \jpkana{イタリアンレストラン}{Itarian resutoran} \jpkana{が}{ga} 
    \jpkana{あります。}{arimasu.}

    \vspace{1.5em}
    \noindent{\Large \color{articolor} \textbf{Arti:} restoran Italia.}
\end{convobox}

\begin{convobox}{34}
    \noindent
    \jpkana{ピザ}{Piza} \jpkana{が}{ga} 
    \jpkana{とても}{totemo} 
    \jpkana{おいしい}{oishii} \jpkana{ので}{node} 
    \jpkana{おすすめ}{osusume} \jpkana{です。}{desu.}

    \vspace{1.5em}
    \noindent{\Large \color{articolor} \textbf{Arti:} Karena pizzanya sangat enak, saya merekomendasikannya.}
\end{convobox}

% --- LESSON BUFFER 4 ---
\begin{lessonbox}{Catatan Belajar: Menyatakan Alasan!}
    Di bagian ini, turisnya memakai kata \textbf{ので (Node)} yang artinya "Karena".
    \begin{itemize}
        \item \textit{Nodo ga kawaita \textbf{node}} = Karena saya haus
        \item \textit{Piza ga oishii \textbf{node}} = Karena pizzanya enak
    \end{itemize}
    Jadi, kalau kamu mau memberi tahu alasan kamu melakukan sesuatu, kamu bisa menempelkan \textbf{"node"} di tengah kalimat. Coba deh bikin kalimat alasanmu sendiri!
\end{lessonbox}

% ================= PERCAKAPAN 35-42 =================
\begin{convobox}{35}
    \noindent
    \jpkana{あっ、}{A,,} 
    \jpkana{どうしよう…}{dou shiyou...}

    \vspace{1.5em}
    \noindent{\Large \color{articolor} \textbf{Arti:} Ah, bagaimana ini...}
\end{convobox}

\begin{convobox}{36}
    \noindent
    \jpkana{スマホ}{Sumaho} \jpkana{の}{no} 
    \jpkana{バッテリー}{batterii} \jpkana{が}{ga} 
    \jpkana{もうすぐ}{mou sugu} 
    \jpkana{なくなります。}{nakunarimasu.}

    \vspace{1.5em}
    \noindent{\Large \color{articolor} \textbf{Arti:} Baterai HP saya sebentar lagi mau habis.}
\end{convobox}

\begin{convobox}{37}
    \noindent
    \jpkana{この}{Kono} 
    \jpword{ちか}{近}{chika}\jpkana{く}{ku} \jpkana{に}{ni} 
    \jpkana{スマホ}{sumaho} \jpkana{を}{o} 
    \jpword{じゅうでん}{充電}{juuden} \jpkana{できる}{dekiru} 
    \jpword{ばしょ}{場所}{basho} \jpkana{は}{wa} 
    \jpkana{ありますか？}{arimasu ka?}

    \vspace{1.5em}
    \noindent{\Large \color{articolor} \textbf{Arti:} Apakah ada tempat untuk bisa mengecas HP di dekat sini?}
\end{convobox}

\begin{convobox}{38}
    \noindent
    \jpkana{コンビニ}{Konbini} \jpkana{で}{de} 
    \jpkana{モバイルバッテリー}{mobairu batterii} \jpkana{が}{ga}

    \vspace{1.5em}
    \noindent{\Large \color{articolor} \textbf{Arti:} Powerbank di minimarket,}
\end{convobox}

\begin{convobox}{39}
    \noindent
    \jpword{か}{買}{ka}\jpkana{えると}{eru to} 
    \jpword{おも}{思}{omo}\jpkana{います。}{imasu.}

    \vspace{1.5em}
    \noindent{\Large \color{articolor} \textbf{Arti:} menurutku bisa dibeli.}
\end{convobox}

\begin{convobox}{40}
    \noindent
    \jpkana{わかりました。}{Wakarimashita.}

    \vspace{1.5em}
    \noindent{\Large \color{articolor} \textbf{Arti:} Saya mengerti.}
\end{convobox}

\begin{convobox}{41}
    \noindent
    \jpkana{いろいろ}{Iroiro} 
    \jpword{おし}{教}{oshi}\jpkana{えてくださって}{ete kudasatte} 
    \jpkana{ありがとうございました。}{arigatou gozaimashita.}

    \vspace{1.5em}
    \noindent{\Large \color{articolor} \textbf{Arti:} Terima kasih banyak sudah mengajariku (memberitahuku) banyak hal.}
\end{convobox}

\begin{convobox}{42}
    \noindent
    \jpkana{それでは、}{Sore dewa,} 
    \jpword{しつれい}{失礼}{shitsurei} \jpkana{します。}{shimasu.}

    \vspace{1.5em}
    \noindent{\Large \color{articolor} \textbf{Arti:} Kalau begitu, saya permisi dulu.}
\end{convobox}

% --- SUMMARY BOX ---
\vspace{2em}
\begin{tcolorbox}[
    colback=blue!5!white,
    colframe=blue!80!black,
    boxrule=3pt,
    arc=5mm,
    title={\huge\bfseries \sffamily 🌟 Rangkuman Bab 2 🌟},
    fonttitle=\color{white},
    attach boxed title to top center={yshift=-4mm},
    boxed title style={colback=blue!80!black, rounded corners},
    left=5mm, right=5mm, top=7mm, bottom=5mm
]
\Large \textbf{Yey, kamu berhasil bantu turis!} Di bab ini kita belajar banyak kosakata super penting kalau jalan-jalan:
\begin{itemize}
    \item \textbf{Arah Jalan:} \textit{Massugu} (Lurus), \textit{Migi / Hidari} (Kanan / Kiri), \textit{Magaru} (Belok).
    \item \textbf{Bertanya Lokasi:} \textit{\dots arimasu ka?} (Apakah ada\dots?)
    \item \textbf{Rekomendasi:} \textit{Osusume} (Rekomendasi / Saran yang bagus).
    \item \textbf{Memberi Alasan:} Gunakan partikel \textit{Node} (Karena).
    \item \textbf{Darurat HP Mati:} \textit{Juuden dekiru} (Bisa mengecas baterai).
\end{itemize}
Bagus sekali! Sekarang kamu sudah siap jadi pemandu wisata dadakan di kotamu. Semangat terus belajarnya ya!
\end{tcolorbox}

\newpage